% PRD submission scaffold for the Λcos paper.
% Step 1 (this file): revtex preamble, title/author/affiliation, abstract.
%                     Body is in body-raw.tex (raw pandoc output, NOT yet cleaned).
% Step 2 will clean up body-raw.tex and \input it here, then drop the TODO block.

\documentclass[
  aps,
  prd,
  reprint,           % two-column PRD format; switch to "onecolumn" if reviewing
  nofootinbib,
  longbibliography
]{revtex4-2}

\usepackage{amsmath}
\usepackage{amssymb}
\usepackage{graphicx}
\usepackage{booktabs}
\usepackage{array}        % for >{} column specifications
\usepackage{calc}         % \real{} for fractional column widths
\usepackage{longtable}    % retained for any longtable that survives cleanup
\usepackage{tabularx}     % X columns auto-wrap long cell content
\usepackage{hyperref}

% Pandoc helpers (see pandoc's default LaTeX template)
\providecommand{\tightlist}{%
  \setlength{\itemsep}{0pt}\setlength{\parskip}{0pt}}

\begin{document}

\title{Apparent Phantom Crossing as Template Bias: \\
       A Non-Phantom Test Case with $\Lambda$cos}

\author{B.~Shatto}
\email{bshatto.pe@gmail.com}
\affiliation{Independent Researcher, St.~Petersburg, FL, USA}

\date{\today}

\begin{abstract}
DESI DR2 reports a 2.8--4.2$\sigma$ preference for dynamical dark energy in combined analyses, with standard $w_0 w_a$ fits favoring an apparent crossing of $w = -1$. We show that a non-phantom expansion history can produce apparent phantom crossings when analyzed with standard two-parameter templates. We construct a specific realization, $\Lambda$cos, and constrain it using Pantheon+ and DESI DR2 BAO. At data-allowed parameter values, the induced CPL distortion is structurally present but modest ($w_0 \approx -1.02$) and of opposite sign in $w_a$, establishing the mechanism without reproducing the DESI best-fit amplitude.

The $\Lambda$cos model yields $H^2(z)/H_0^2 = \alpha(1+z)^3 - \beta(1+z) + \Omega_\Lambda$, where $\alpha$ and $\beta$ are determined by a single parameter $s_0$ and a fixed reference value $\Omega_\Lambda = 0.685$. Under the fiducial-matter diagnostic split, the model satisfies $w_{\rm eff}(z) > -1$ at all redshifts, recovers the fiducial flat $\Lambda$CDM limit as $s_0 \to 0$, and contains a negative $(1+z)^1$ term in $H^2$ absent from the four canonical FLRW density scalings. A joint MCMC fit to Pantheon+ (1701 SNe Ia) and DESI DR2 BAO (13 data points) yields $\Delta\chi^2 = +0.11$ relative to flat $\Lambda$CDM at the same parameter count, with $s_0 < 0.18$ (95\% CL, flat prior). The constraint is robust across prior choices ($s_0 < 0.12$--$0.21$) and $\Omega_\Lambda$ variations ($0.68$--$0.715$). Adding compressed Planck distance priors shifts the preferred $\Omega_\Lambda$ in both $\Lambda$CDM and $\Lambda$cos; allowing $\Omega_\Lambda$ to vary gives statistically equivalent fits ($\Delta\chi^2 = +0.28$ at equal parameter count).
\end{abstract}

\maketitle

% Step 2 done (cleanup_body.py): structural fixes, heading-level bump,
% math-mode underscores, > / < conversion. body.tex is the cleaned body.
%
% Remaining for later steps:
%   Step 3: numeric citations {[}N{]} -> \cite{...}, write references.bib
%   Step 4: [Figure N: caption] -> \begin{figure}...\end{figure} environments
%           pointing at /Users/blake/lambda-cos/figures/figN_*.pdf
%   Step 5: § cross-refs in text -> Sec.~\ref{sec:...} after labels are added

\input{body}

\bibliographystyle{apsrev4-2}
\bibliography{references}

\end{document}
